\subsection{Effective computation at each fixed time}
\label{profile:time}

For rational $T>0$, the preceding pruning construction computes both
$G(T)$ and $H(T)$ with any prescribed rational accuracy. Its exterior
potential event is independent of the future tagged potential
genealogy, and both actual genealogies are subsets of that genealogy.
The potential total core count has mean $e^t$. Consequently, uniformly
for $0\le t\le T$,
\begin{equation}\label{profile:time-error}
 |F_x(t)-F_{\pi_{R,M}x}(t)|
 \le \varepsilon_{R,M}(T)
 :=e^T\left(\frac{T^R}{R!}+\frac{e^T-1}{M}\right).
\end{equation}
If $\mathcal P_{R,M}$ is the finite set of pruned profiles, then
\begin{equation}\label{profile:time-max}
 \max_{P\in\mathcal P_{R,M}}F_P(t)
 \le G(t)\le
 \max_{P\in\mathcal P_{R,M}}F_P(t)+\varepsilon_{R,M}(T).
\end{equation}
The first inequality uses the admissibility of each profile, proved
by finite inert-boundary completion in Section~\ref{model:section}.
Choose $R,M$ effectively to make the displayed error arbitrarily small.

For a fixed profile with $n$ real vertices, enumerate the finite state
space with real population at most $N$, including genealogical marks,
and kill on the next birth. All transition rates are rational. If
$Y_n(t)$ is a Yule process started at $n$, the omitted mean is at most
\begin{equation}\label{profile:population-error}
 \mathbb E[Y_n(t);Y_n(t)>N]
 \le\frac{\mathbb EY_n(t)^2}{N}
 =\frac{n(n+1)e^{2t}-ne^t}{N}.
\end{equation}
Indeed, the real birth rate is at most the population; inert choices
are null. Couple with $Y_n$ so that a crossing of the retained cutoff
forces $Y_n(t)>N$, and use $C_{\rm tag}(t)\le Y_n(t)$.

A rational uniformization rate exceeding every finite-state exit rate
expresses its matrix exponential as a Poisson-weighted series of
powers of a rational substochastic matrix. Its partial sums and a
certified Poisson tail give rational enclosures. Elementary rational
enclosures for exponentials make \eqref{profile:population-error}
effective as well.

Finally every $F_x$ is $e^T$-Lipschitz on $[0,T]$: its core-birth
intensity is at most $F_x(t)\le e^t$. The supremum $G$ has the same
Lipschitz constant. Rational-grid quadrature, combined with
\eqref{profile:time-max} and the population cutoff, computes
$H(T)=1+\int_0^TG(t)\,dt$. This argument concerns a fixed time and
uses no convergence rate for $t^{-1}\log G(t)$.
